\begin{abstract}
% ── Instructions ─────────────────────────────────────────────────────────
% IEEE conference: 150–250 words, single paragraph, no citations, no maths.
% Cover: motivation, approach, key results, conclusions.
% ─────────────────────────────────────────────────────────────────────────

SLAM systems accumulate drift over long trajectories, degrading trajectory
accuracy and map consistency---particularly for 2D LiDAR systems in
featureless indoor environments. Infrastructure-assisted approaches typically
register fixed cameras onto a pre-built SLAM map, so any drift in that map
biases subsequent camera poses. This paper presents a method that jointly
estimates the environment map and the poses of ceiling-mounted cameras,
enabling cameras to act as global anchors without a separate registration
stage. When a camera observes the robot, a visual constraint is inserted into
the pose graph, tying the robot trajectory to the camera's map-frame location.
Unlike loop closure alone, these repeated absolute constraints suppress drift
in visually ambiguous and revisited areas. The approach integrates into a
graph-based SLAM backend via pose-graph optimisation with no modifications
to the core pipeline. Experiments in simulated indoor environments demonstrate
improved map consistency and reduced absolute trajectory error compared with a
LiDAR-only baseline, at modest additional computational cost.

\end{abstract}

\begin{IEEEkeywords}
SLAM, visual constraints, pose-graph optimisation, fixed cameras
\end{IEEEkeywords}
